\documentclass[11pt]{article}
\usepackage[margin=1in]{geometry}
\usepackage{amsmath,amssymb,amsthm}
\usepackage[colorlinks=true,urlcolor=blue,linkcolor=blue]{hyperref}
\newtheorem{theorem}{Theorem}
\newtheorem{lemma}{Lemma}
\newcommand{\G}{\Gamma}
\newcommand{\dd}{\,d}
\title{Coefficientwise log-concavity of the normalized\\
incomplete beta function\\\large MathDB 331384}
\author{Principia Math In Progress}
\date{11 September 2026}
\begin{document}
\maketitle

\section*{Task 1. Resolve the first-parameter conjecture}

For $a,b>0$, let $I_x(a,b)$ be the regularized incomplete beta function and set
\[
 f_b(a;x)=x^{-a}(1-x)^{-b}I_x(a,b)
 =\frac1{\G(b)}\sum_{n=0}^\infty
       \frac{\G(a+b+n)}{\G(a+1+n)}x^n,\qquad |x|<1.
\]
The series identity follows from Euler's transformation of the standard
hypergeometric representation of $I_x$; it also identifies the formal
coefficients under consideration. Write
\[
 f_b(a+\alpha;x)f_b(a+\beta;x)
 -f_b(a;x)f_b(a+\alpha+\beta;x)
 =\sum_{m=0}^\infty \psi_m x^m .
\]
The following proves Conjecture 2 of Karp~\cite{K}, including the
unrestricted positive shifts in the previously unsettled $b>1$ case.

\begin{theorem}
For every $a,\alpha,\beta>0$ and every integer $m\geq0$,
\[
 \psi_m>0\quad(b>1),\qquad
 \psi_m<0\quad(0<b<1),\qquad
 \psi_m=0\quad(b=1).
\]
Thus $a\mapsto f_b(a;x)$ is strictly coefficientwise Wright
log-concave for $b>1$ and strictly coefficientwise Wright log-convex
for $0<b<1$.
\end{theorem}

The $0<b<1$ assertion was already proved in the source's Corollary 1.
We give its short moment proof below to make the result self-contained.
The main step for $b>1$ is a finite-sum inequality.

\section*{A finite-sum inequality}

Put
\[
 g_b(t)=\frac{\G(t+b)}{\G(t+1)},\qquad t>0.
\]

\begin{lemma}\label{lem:kernel}
Let $b>1$, $u>v>0$ and $m\geq0$ be an integer. Then
\begin{equation}
 \sum_{k=0}^m g_b(u+k)g_b(v+m-k)
       \bigl(q^{u+k}-q^{v+m-k}\bigr)<0
       \qquad(0<q<1).
 \label{eq:kernel}
\end{equation}
\end{lemma}

\begin{proof}
Set $T=u+v+m$, $w(t)=g_b(t)g_b(T-t)$ and
\[
 P_u(q)=\sum_{k=0}^m w(u+k)q^{u+k},\qquad
 P_v(q)=\sum_{k=0}^m w(v+k)q^{v+k},\qquad
 F=P_u-P_v.
\]
Since $w(t)=w(T-t)$, $F$ is the left side of~\eqref{eq:kernel}.
Let $\theta=q\,d/dq$ and introduce the differential operator
\[
 \mathcal L=\theta(T+b-\theta)-q(\theta+b)(T-\theta).
\]
The weight recurrence
\begin{equation}
 (t+1)(T-t+b-1)w(t+1)=(t+b)(T-t)w(t)
 \label{eq:weights}
\end{equation}
shows that all interior terms in $\mathcal LP_u$ cancel.
Subtracting the two boundary expressions gives
\begin{equation}
 \mathcal LF=R(q):=
 C\bigl(q^u+q^{T-u+1}\bigr)
 -D\bigl(q^v+q^{T-v+1}\bigr),
 \label{eq:forcing}
\end{equation}
where
\[
 C=u(T+b-u)w(u),\qquad D=v(T+b-v)w(v).
\]

We first determine the sign changes of $R$. Define
$r_b(t)=\G(t+b)/\G(t)$. Its logarithm is strictly concave, since
\[
 (\log r_b)''(t)=\psi'(t+b)-\psi'(t)<0;
\]
the last inequality follows, for example, from
$\psi'(t)=\sum_{j\geq0}(t+j)^{-2}$.
Moreover,
\[
 C=r_b(u)r_b(T+1-u),\qquad
 D=r_b(v)r_b(T+1-v).
\]
The two arguments in either product have sum $T+1$, and
\[
 \left|u-\frac{T+1}{2}\right|
 <\left|v-\frac{T+1}{2}\right|,
\]
because $u>v$ and $u+v=T-m<T+1$.
Strict log-concavity therefore gives $C>D>0$.

The ratio
\[
 \frac{q^u+q^{T-u+1}}{q^v+q^{T-v+1}}
 =\frac{\cosh(\rho s)}{\cosh(\sigma s)},\qquad
 s=-\log q,\quad
 \rho=\left|u-\frac{T+1}{2}\right|,\quad
 \sigma=\left|v-\frac{T+1}{2}\right|,
\]
increases strictly from $0$ to $1$ as $q$ increases from $0$ to $1$.
Indeed $0\leq\rho<\sigma$, and its logarithmic derivative in $s>0$
is $\rho\tanh(\rho s)-\sigma\tanh(\sigma s)<0$.
Thus $R$ has exactly one zero in $(0,1)$, with signs $-,+$.

We next construct a positive solution of the homogeneous equation.
Let
\begin{equation}
 Y(q)=q^{T+b}h(q),\qquad
 h(q)=\sum_{j=0}^\infty
 \frac{(b)_j(T+2b)_j}{(T+b+1)_j\,j!}q^j .
 \label{eq:ground}
\end{equation}
The coefficients are positive and the series converges for $0<q<1$.
Its coefficient recurrence gives $\mathcal LY=0$ directly.
In particular $Y(q)>0$ and $Y(q)\geq q^{T+b}$.

For use at the endpoints, the beta integral gives
\begin{equation}
 h(q)=\frac1{B(b,T+1)}
 \int_0^1\frac{t^{b-1}(1-t)^T}{(1-qt)^{T+2b}}\dd t .
 \label{eq:euler}
\end{equation}
Near zero, $Y=O(q^{T+b})$ and $Y'=O(q^{T+b-1})$.
Near one, writing $\varepsilon=1-q$ and using $q\geq1/2$,
$t^{b-1}\leq1$ and $1-qt\geq(\varepsilon+1-t)/2$, we obtain
\[
 h(q)=O(\varepsilon^{1-2b}),\qquad
 h'(q)=O(\varepsilon^{-2b}).
\]
For example, the first estimate follows by bounding the integral by
a constant times
$\int_0^1 z^T(\varepsilon+z)^{-T-2b}\dd z$
and substituting $z=\varepsilon y$; the resulting integral on
$(0,\infty)$ is finite because $b>1$.
Differentiating~\eqref{eq:euler} proves the second estimate in the same
way. Consequently
\begin{equation}
 Y=O((1-q)^{1-2b}),\qquad
 Y'=O((1-q)^{-2b}) \quad(q\uparrow1).
 \label{eq:endpoint}
\end{equation}

In ordinary derivatives, equation~\eqref{eq:forcing} is
\[
 -q^2(1-q)F''
 +q\bigl(T+b-1-q(T-b-1)\bigr)F'
 -qbTF=R.
\]
Set
\[
 p(q)=q^{1-T-b}(1-q)^{2b},\qquad
 W(q)=p(q)\bigl(Y(q)F'(q)-Y'(q)F(q)\bigr).
\]
The homogeneous equation for $Y$ and the displayed equation for $F$
give the exact identity
\begin{equation}
 W'(q)=-\frac{p(q)Y(q)}{q^2(1-q)}R(q).
 \label{eq:wronskian}
\end{equation}
Both endpoint limits of $W$ are zero. At zero,
$F=O(q^v)$, $F'=O(q^{v-1})$, so $W=O(q^v)$.
At one, the finite sums give $F(1)=0$, $F=O(1-q)$ and $F'=O(1)$;
using~\eqref{eq:endpoint} gives $W=O(1-q)$.

Since $R$ has signs $-,+$ and all other factors in
\eqref{eq:wronskian} are positive, $W$ first increases from zero
and then decreases to zero. Hence $W(q)>0$ throughout $(0,1)$.
It follows that
\[
 \left(\frac F Y\right)'=\frac{W}{pY^2}>0.
\]
Finally $F/Y\to0$ as $q\uparrow1$, since $F\to0$ and
$Y\geq q^{T+b}$. Thus $F/Y<0$, and $F<0$, as required.
\end{proof}

\section*{Proof of coefficientwise log-concavity}

For fixed $m$ define the symmetric function
\[
 H_m(u,v)=\sum_{k=0}^m g_b(u+k)g_b(v+m-k).
\]
Suppose $b>1$ and hold $u+v$ fixed. Differentiating with respect to $u$
(so $dv/du=-1$) and using the digamma integral gives
\begin{align}
 \frac{d}{du}H_m(u,v)
 &=\int_0^\infty
       \frac{e^{-s}-e^{-bs}}{1-e^{-s}}
       \sum_{k=0}^m g_b(u+k)g_b(v+m-k)
       \bigl(e^{-(u+k)s}-e^{-(v+m-k)s}\bigr)\dd s .
 \label{eq:derivative}
\end{align}
Indeed
\[
 \frac{d}{dt}\log g_b(t)
 =\psi(t+b)-\psi(t+1)
 =\int_0^\infty e^{-ts}
       \frac{e^{-s}-e^{-bs}}{1-e^{-s}}\dd s .
\]
There are only finitely many terms, and the displayed integrals converge
absolutely. For $u>v$, the finite sum in~\eqref{eq:derivative} is
strictly negative by Lemma~\ref{lem:kernel}, with $q=e^{-s}$.
The remaining factor is strictly positive. Thus
$H_m(u,v)$ strictly decreases as $|u-v|$ increases at fixed $u+v$.

The pairs $(a+\alpha,a+\beta)$ and $(a,a+\alpha+\beta)$ have the
same sum, whereas $|\alpha-\beta|<\alpha+\beta$. Therefore
\[
 \psi_m=\frac{H_m(a+\alpha,a+\beta)-H_m(a,a+\alpha+\beta)}
                {\G(b)^2}>0
 \qquad(b>1).
\]
This applies also to $m=0$ and requires no integrality of either shift.

\section*{The other values of \(b\)}

For $0<b<1$, the beta integral gives a positive measure
\[
 \dd\nu_b(s)=
 \frac{s^{b-1}(1-s)^{-b}}{\G(1-b)}\dd s,\qquad 0<s<1,
 \qquad g_b(t)=\int_0^1s^t\dd\nu_b(s).
\]
Put $K_m(s,t)=\sum_{k=0}^m s^kt^{m-k}$, which is positive and symmetric
on $(0,1)^2$. Multiplying the moment integrals and symmetrizing yields
\[
 \psi_m=-\frac1{2\G(b)^2}
 \int_0^1\int_0^1
 (st)^a(s^\alpha-t^\alpha)(s^\beta-t^\beta)
 K_m(s,t)\dd\nu_b(s)\dd\nu_b(t)<0.
\]
The integrand apart from its leading minus sign is positive off the
diagonal, a set of full product measure. All integrals are finite.
At $b=1$, $g_1(t)=1$ and $f_1(a;x)=1/(1-x)$ is independent of $a$,
so all the coefficients vanish. This completes the theorem.

\section*{Verification and source status}

The proof is analytic and does not depend on numerical sign tests.
The accompanying verifier checks the finite-sum differential identity
with exact rational arithmetic, the homogeneous-series recurrence,
and independent numerical instances of the integral identities.
The live MathDB record, published statement and public claim sources
were checked on 11 September 2026; no matching full solution claim was
located. The dated records and search limits are in
\texttt{literature-audit.md}.

\section*{Tasks 2--5}

\begin{enumerate}
\setcounter{enumi}{1}
\item Determine which nonnegative sequences $(c_n)$ preserve the result
when $g_b(a+n)$ is replaced by $c_n g_b(a+n)$, including the full class
of log-concave sequences without internal zeros.
\item For $b>1$, derive explicit positive lower bounds and matching upper
bounds for each coefficient $\psi_m$, in terms of the parameters and its order.
\item Prove or refute the analogous unrestricted coefficientwise
log-concavity conjecture for shifts in the second parameter of $I_x(a,b)$.
\item Establish a \(q\)-gamma analogue of the finite-sum inequality
in Lemma~\ref{lem:kernel} and determine the corresponding parameter ranges.
\end{enumerate}

\begin{thebibliography}{9}
\bibitem{K}
D.~B. Karp, \emph{Normalized incomplete beta function:
log-concavity in parameters and other properties},
Journal of Mathematical Sciences \textbf{217} (2016), 91--107,
Conjecture 2, p.~99, and Corollary 1, p.~101.
\href{https://doi.org/10.1007/s10958-016-2958-z}
{doi:10.1007/s10958-016-2958-z}.
\href{https://dmkrp.wordpress.com/wp-content/uploads/2016/07/karp2016.pdf}
{Published full text};
\href{https://arxiv.org/abs/1509.05120}{arXiv:1509.05120}.
\end{thebibliography}
\end{document}
